% =========================================================
%  frontmatter/preface.tex
%  Kata pengantar
% =========================================================

\chapter*{Kata Pengantar}
\addcontentsline{toc}{chapter}{Kata Pengantar}
\thispagestyle{plain}

Buku ajar ini disusun untuk mendukung perkuliahan \NamaMK\ (\KodeMK) yang
berbobot \BobotSKS\ pada Program Studi \NamaProdi, \NamaFakultas,
\NamaPT. Penyusunannya berangkat dari satu pertanyaan sederhana: bagaimana
menyajikan materi pemrograman permainan agar mahasiswa benar-benar mampu
membangun sesuatu, bukan sekadar mengenal istilah.

Mata kuliah ini dirancang dengan pendekatan \istilah{Outcome-Based
Education} (OBE). Konsekuensinya, seluruh materi, aktivitas, dan asesmen
dalam buku ini diarahkan pada capaian pembelajaran yang telah ditetapkan
lebih dahulu. Setiap bab menyebut Sub-CPMK yang dicapainya, menyediakan
praktikum yang menghasilkan artefak nyata, dan menutup dengan checklist
kompetensi serta instrumen asesmen. Dengan demikian pembaca dapat menelusuri
secara eksplisit hubungan antara apa yang dipelajari, apa yang dikerjakan,
dan apa yang dinilai.

Pilihan teknologi pada buku ini mengikuti urutan yang disengaja. Materi
pertama sampai kesembilan dikerjakan dengan Python dan Pygame tanpa bantuan
\istilah{game engine}. Mahasiswa menulis sendiri \istilah{game loop},
menghitung sendiri gerak berbasis \istilah{delta time}, menerapkan sendiri
hukum fisika sederhana, dan mengimplementasikan sendiri algoritma
kecerdasan buatan untuk NPC, termasuk \istilah{Finite State Machine} dan
pencarian jalur A*. Cara ini membuat mekanisme internal permainan terlihat
jelas dan dapat diuji. Setelah fondasi tersebut terbentuk, materi kesepuluh
sampai keempat belas memindahkan pekerjaan ke Unity, tempat mahasiswa
belajar memanfaatkan \istilah{game engine} untuk membangun level, mengimpor
aset, menata antarmuka, dan mengintegrasikan audio. Perbandingan antara
kedua pendekatan ini menjadi salah satu bahan refleksi penting pada akhir
semester.

Buku ini tidak dimaksudkan sebagai pengganti dokumentasi resmi. Python,
Pygame, dan Unity memiliki dokumentasi yang lengkap dan terus diperbarui;
buku ini berperan sebagai pemandu yang mengurutkan, menjelaskan, dan
mengaitkan materi tersebut dengan kebutuhan mata kuliah. Rujukan
\mbox{[R1]}--\mbox{[R14]} pada daftar pustaka adalah sumber utama yang
seyogianya dibaca mahasiswa, sedangkan \mbox{[R15]}--\mbox{[R35]} memuat
dokumen kurikulum, bahan ajar, dan studi kasus dari berbagai perguruan
tinggi Indonesia yang menjadi bahan pembanding penyusunan buku ini.

Umpan balik dari dosen pengampu dan mahasiswa sangat berharga bagi
perbaikan edisi berikutnya. Saran dapat disampaikan langsung kepada penulis
melalui Program Studi \NamaProdi, \NamaPT.

\vspace{0.8cm}

\begin{flushright}
Bandung, \BookYear\\
\BookAuthor
\end{flushright}

\clearpage
